\section{Method Details}
\label{app:method_details}

\subsection{Diagnostics}
\label{app:diagnostics}

Three diagnostic procedures complement the primary DIF analysis.

\paragraph{Local independence (Yen's Q3).} The MH procedure assumes that item responses are conditionally independent given $\theta$. Yen's Q3 statistic measures residual correlations between item pairs after conditioning on the latent trait. We compute Q3 for all within-subject and cross-subject item pairs, with $|Q3| > 0.20$ as the threshold for local dependence. Elevated within-subject dependence would indicate that items within the same subject share variance beyond what $\theta$ captures, which could inflate within-subject DIF rates.

\paragraph{Domain-specificity.} We test whether temporal DIF concentrates in domains where models improved most. At the ecological level, we compute Spearman $\rho$ between subject-level C\% and subject-level accuracy improvement from 2023 to 2024. At the item level, logistic regression predicts item-level DIF~C classification from domain-level accuracy change, with and without item property controls (difficulty, discrimination). This decomposition quantifies how much of the observed instability is attributable to aggregate domain-level gains versus item-specific behavioral change.

\paragraph{Instability decomposition.} When domain accuracy improvement is included alongside item-level controls, we compute the counterfactual C\% that would remain after removing domain-level effects. The gap between observed and counterfactual C\% measures the contribution of domain-level improvement to the total instability.

\subsection{Statistical Considerations}
\label{app:statistical}

\paragraph{Multiple testing.} The ETS~C classification imposes a dual threshold (effect size $|\Delta_{\mathrm{MH}}| \geq 1.5$ and significance $p < 0.05$), which provides implicit control over false positives. As a robustness check, we apply Benjamini-Hochberg FDR correction at $q = 0.05$. The random-split sanity check provides an empirical estimate of the false positive rate under the null hypothesis.

\paragraph{Threshold sensitivity.} We report results across a range of DIF thresholds ($|\Delta_{\mathrm{MH}}| \in \{1.0, 1.5, 2.0, 2.5\}$) to verify that qualitative conclusions do not depend on the specific ETS~C cutoff. The standard threshold of $|\Delta_{\mathrm{MH}}| \geq 1.5$ corresponds to the ETS~C definition used throughout the paper.

\paragraph{Architecture heterogeneity.} The METABENCH MMLU population consists of 99.98\% decoder-only models (5,226 of 5,227). This near-complete architectural homogeneity eliminates architecture as a confound for DIF analysis. We verify this by excluding the single encoder-decoder model and confirming identical results.

\section{Structural Analysis Details}
\label{app:structural}

Three progressively refined approaches fail to reduce the temporal false positive rate below 31\%, revealing a structural property of the data rather than a methodological limitation. Matched-ability quintile analysis restricts comparisons to models of similar total score and yields enrichment ratios of 0.94--1.04, with temporal FPR unchanged at approximately 0.31. IRT $\theta$-conditioning with 20 ability strata achieves $r(\theta, \mathrm{total\ score}) = 0.9986$ and FPR $= 0.312$, confirming that the external total score already captures nearly all ability variance relevant to DIF computation. Within-subject DIF analysis on 5 representative subjects, intended to isolate subject-specific effects, finds mean C\% $= 38.1\%$, exceeding the 25\% threshold that would indicate viable bifactor structure. \Cref{tab:structural_comparison} summarizes these results.

The temporal FPR reflects genuine item-level instability that persists regardless of how ability is measured or conditioned. Yen's Q3 diagnostic reveals mild local dependence within subjects (11.3\% of within-subject item pairs exceed $|Q3| > 0.20$, vs.\ 2.2\% cross-subject; $5{\times}$ ratio, $p = 2.3 \times 10^{-9}$), which contributes to elevated within-subject DIF rates but is not the primary driver: the cross-subject matching protocol addresses this bias, and temporal C\% persists under external matching (full Q3 distributions in Appendix~Figure~A2).

\begin{table}[h]
\centering
\caption{Structural approaches to reducing temporal DIF. None reduces baseline FPR below 31\%, indicating that the observed instability is not an artifact of ability matching. Subscript $\pm$ values denote standard deviations across 5 injection seeds (baseline row only).}
\label{tab:structural_comparison}
\small
\begin{tabular}{lcccc}
\toprule
Approach & Mechanism & Temporal FPR & TPR ($p{=}0.5$) & Outcome \\
\midrule
Baseline (external) & Cross-benchmark total score & $0.314_{\pm.009}$ & $0.834_{\pm.027}$ & Reference \\
Matched-ability & Within-quintile temporal DIF & ${\sim}$0.31 & -- & Enrich.\ 0.94--1.04 \\
$\theta$-conditioning & $\theta$ replaces total score & 0.312 & 0.876 & $r(\theta,\mathrm{total}) {=} 0.9986$ \\
Within-subject & Per-subject DIF & -- & -- & C\% $= 38.1\%$ ($>$25\%) \\
\bottomrule
\end{tabular}
\end{table}

\section{Additional Figures}
\label{app:figures}

\begin{figure}[h]
\centering
\includegraphics[width=\columnwidth]{fig_3_within_subject_matching.pdf}
\caption{Cross-subject vs.\ within-subject matching for semi-synthetic DIF detection. Cross-subject matching eliminates FPR inflation ($\Delta\mathrm{FPR} = 0.000$) while preserving detection power ($\Delta\mathrm{TPR} > 0.3$ for 4/5 subjects at $p = 0.5$). Within-subject matching inflates FPR by ${+}0.186$ on average.}
\label{app:fig_matching}
\end{figure}

\begin{figure}[h]
\centering
\includegraphics[width=\columnwidth]{fig_4_downstream_impact.pdf}
\caption{Downstream impact of DIF-based item filtering on model rankings. Despite $\rho \approx 0.997$ overall, individual model families show systematic rank shifts when evaluated on stable items only.}
\label{app:fig_downstream}
\end{figure}

\section{Cross-Benchmark Comparison}
\label{app:gsm8k_details}

\begin{table}[h]
\centering
\caption{Cross-benchmark comparison of DIF analysis on MMLU and GSM8K. Both benchmarks exhibit temporal item instability with validated semi-synthetic detection power. Standard and external matching use different exploitation rates ($p$) because each benchmark's optimal operating point differs; external matching results are compared at $p = 0.5$ for both. Subscript $\pm$ values denote standard deviations across 5 random injection seeds.}
\label{tab:gsm8k_comparison}
\small
\begin{tabular}{lcc}
\toprule
Metric & MMLU & GSM8K \\
\midrule
Items & 12,508 & 1,319 \\
Models & 5,227 & 6,702 \\
Temporal C\% & 31.3\% & 18.9\% \\
Sanity C\% & 0.18\% & 0.3\% \\
$\Delta$TPR ($p{=}0.7$, standard) & -- & $+0.818_{\pm.006}$ \\
$\Delta$TPR ($p{=}0.5$, external) & $+0.508_{\pm.027}$ & $+0.514_{\pm.028}$ \\
\bottomrule
\end{tabular}
\end{table}

\section{Limitations}
\label{app:limitations}

\paragraph{(a) Ecological vs.\ item-level decomposition.} The ecological correlation between subject-level C\% and accuracy improvement is $\rho = 0.307$ ($p = 0.02$), but item-level decomposition shows that domain improvement explains only 2.4 percentage points (7.6\%) of 31.3\% C\%. The remaining 92.4\% cannot be fully attributed to contamination; other sources of capability change (training methodology, data composition, optimization procedures) may contribute, and our analysis does not isolate these factors.

\paragraph{(b) Instability attribution.} The 31\% C\% measures items that function differently across cohorts, but the method does not distinguish contamination from legitimate capability change. DIF identifies \emph{that} items changed, not \emph{why}. The semi-synthetic framework validates detection power for one specific mechanism (differential exposure), but real-world instability likely reflects multiple interacting causes.

\paragraph{(c) Semi-synthetic injection fidelity.} The injection procedure flips incorrect responses to correct at a specified rate, simulating an idealized contamination scenario. Real contamination may manifest as partial knowledge, format memorization, or selective retrieval rather than wholesale answer flipping. The gap between our injection model and realistic contamination mechanisms is unknown.

\paragraph{(d) Benchmark scope.} The primary analysis covers MMLU with cross-benchmark validation on GSM8K (\Cref{sec:gsm8k}). GSM8K replication confirms framework portability and provides the GSM1K contamination correlation ($N = 33$), but full temporal DIF characterization on diverse benchmarks (code generation, long-form reasoning, multimodal tasks) has not been conducted. The GSM1K sample ($N = 33$) produces wide confidence intervals that limit the precision of the negative correlation estimate.

\paragraph{(e) Cross-subject baseline false positive rate.} The cross-subject matching protocol achieves $\Delta\mathrm{FPR} = 0.000$ for injected contamination but does not reduce the baseline temporal FPR of 33--53\% across subjects. This rate is stable across injection levels, confirming it reflects genuine temporal instability rather than methodological error, but it means that a large fraction of flagged items in temporal comparisons are not attributable to the specific contamination mechanism under test.

\paragraph{(f) RL model robustness.} Reinforcement learning from human feedback may alter item-level response patterns in ways that differ from pretraining-based contamination. The theoretical motivation for DIF robustness to RL (RL modifies the ability distribution rather than item parameters) is plausible but empirically unvalidated with adequate sample size: METABENCH contains only $N = 13$ models with documented reasoning-RL training.

\paragraph{(g) Model family coverage.} Downstream impact analysis is limited to families with $N \geq 20$ models (Llama, Mistral, Qwen, Phi). Among the top 50 rank movers, 49 of 50 belong to the ``unknown'' family category in METABENCH, precluding interpretation of what training choices produce large rank shifts.

\paragraph{(h) Future directions.} Cross-benchmark generalization to code, multimodal, and long-form reasoning benchmarks would test DIF applicability beyond factual knowledge tasks. RL model cohorts with adequate $N$ would enable empirical validation of the RL robustness hypothesis. MIMIC models and propensity-score matching could address the ability confounding that the current MH-DIF framework does not fully resolve.
